% !TEX program = xelatex
% Minimal document using the arabdoc package.
% The \documentclass{arabdoc-article} line replaces 15-25 lines of
% manual preamble (fontspec, polyglossia, bidi, font setup,
% hyperref, geometry, section formatting, etc.) with a single call.
%
% Compare this to the basic-arabic example which requires:
%   \usepackage{fontspec}
%   \usepackage{polyglossia}
%   \setmainlanguage[...]{arabic}
%   \setotherlanguage{english}
%   \newfontfamily\arabicfont[...]{Amiri}
%   \usepackage{bidi}
%   \usepackage{geometry}
%   \usepackage{hyperref}
%   ... and more
%
% With arabdoc, all of the above is handled automatically.
\documentclass[12pt, a4paper, amiri, captions]{arabdoc-article}

\begin{document}

\title{مستندٌ عربيٌّ بحزمة arabdoc}
\author{د. نبراس أبو الذهب}
\date{2025}
\maketitle

\section{المقدمة}

تُتيحُ حزمةُ \LR{arabdoc} إعدادَ مستندٍ عربيٍّ كاملٍ بسطرٍ واحدٍ
من الإعداد. بدلاً من كتابةِ خمسةَ عشرَ إلى خمسةٍ وعشرينَ سطراً
في المُقدّمةِ لتحميلِ الحزمِ وتحديدِ الخطوطِ واللغةِ والاتجاه،
يَكفي استخدامُ \LR{\textbackslash documentclass\{arabdoc-article\}}.

\section{الخيارات المتاحة}

يَدعمُ \LR{arabdoc} خياراتٍ متعددةً:
الخطّ (\LR{amiri}، \LR{scheherazade}، \LR{lateef})،
الأرقام (\LR{maghrib}، \LR{arabicdigits})،
التقويم (\LR{gregorian}، \LR{islamic})،
والتسمياتِ العربيةَ (\LR{captions}).

\section{الخاتمة}

هذا المثالُ يُوضّحُ بساطةَ الإعدادِ معَ \LR{arabdoc}، حيثُ
يَتولّى الحزمةُ كلَّ تفاصيلِ التنضيدِ العربيّ تلقائياً.

\end{document}
